% !TeX root = ../presentacion.tex
%=======================================================================
%  config/tema.tex — identidad visual de la presentación
%
%  Base: metropolis. Encima, la paleta del informe (E1), para que la
%  charla y el documento se lean como una sola entrega y no como dos
%  trabajos distintos.
%=======================================================================

\usetheme{metropolis}

\metroset{
  % Sin lámina separadora entre bloques: en una charla de siete minutos, once
  % separadoras son casi un minuto de pantallas que no dicen nada. Las
  % secciones se conservan para el índice del PDF y la barra de progreso.
  sectionpage  = none,
  subsectionpage = none,
  numbering    = fraction,     % «12 / 34» abajo a la derecha
  progressbar  = frametitle,   % barra fina bajo el título de cada lámina
  block        = fill,         % los bloques llevan fondo, no solo título
  background   = light
}

%------------------------------ Paleta ---------------------------------
% acento: el mismo azul del informe (es también el azul de Structurizr, con
% lo que los diagramas C4 no desentonan al proyectarse).
\definecolor{acento}{HTML}{1168BD}
\definecolor{acentoOscuro}{HTML}{0B4C8C}
\definecolor{acentoClaro}{HTML}{E7F0FA}
\definecolor{texto}{HTML}{23272A}
\definecolor{textoSuave}{HTML}{5B6570}
\definecolor{fondo}{HTML}{FFFFFF}
\definecolor{marcador}{HTML}{B45309}   % ámbar: los [Insertar ...] pendientes
\definecolor{ok}{HTML}{0F766E}         % verde azulado: lo que sí se cumple
\definecolor{alerta}{HTML}{B3261E}     % rojo: lo que no, o lo que se cayó
\definecolor{guiaback}{HTML}{F3F4F6}
\definecolor{guiaframe}{HTML}{9AA3AF}

%--------------------- Tintas del mapa de arquitectura -----------------
% Versiones claras de los colores que Structurizr da a cada tipo de
% elemento, para que la lámina de arquitectura y los diagramas C4 del
% respaldo se lean como la misma cosa. Claras a propósito: el texto va
% encima en negro y tiene que leerse desde el fondo del aula.
\definecolor{tintaEdge}{HTML}{E4E7EA}
\definecolor{tintaFront}{HTML}{EDE3F7}
\definecolor{tintaGw}{HTML}{FAE7D0}
\definecolor{tintaMs}{HTML}{DDEFE2}
\definecolor{tintaDb}{HTML}{EADCF3}

\setbeamercolor{normal text}{fg=texto, bg=fondo}
\setbeamercolor{alerted text}{fg=acento}
\setbeamercolor{example text}{fg=ok}
\setbeamercolor{frametitle}{fg=fondo, bg=acento}
\setbeamercolor{progress bar}{fg=acentoOscuro, bg=acentoClaro}
\setbeamercolor{title separator}{fg=acento}
\setbeamercolor{progress bar in head/foot}{fg=acento}
\setbeamercolor{progress bar in section page}{fg=acento}
\setbeamercolor{block title}{fg=acentoOscuro, bg=acentoClaro}
\setbeamercolor{block body}{fg=texto, bg=guiaback}
\setbeamercolor{itemize item}{fg=acento}
\setbeamercolor{itemize subitem}{fg=textoSuave}
\setbeamercolor{section in toc}{fg=texto}

%--------------------------- Tipografía --------------------------------
\setbeamerfont{frametitle}{size=\large, series=\bfseries}
\setbeamerfont{title}{size=\LARGE, series=\bfseries}
\setbeamerfont{subtitle}{size=\normalsize, series=\mdseries}
\setbeamerfont{caption}{size=\scriptsize}
\setbeamerfont{footnote}{size=\tiny}

% Viñetas discretas: en pantalla, un triangulito relleno lee mejor que el
% guion largo por defecto de metropolis.
\setbeamertemplate{itemize item}{\raisebox{0.15ex}{\scriptsize\ding{110}}}
\setbeamertemplate{itemize subitem}{\raisebox{0.15ex}{\tiny\ding{110}}}

%--------------------------- Pie de página -----------------------------
% Metropolis solo pone el número. Se le añade el nombre del curso a la
% izquierda: sirve de marca de agua discreta si alguien fotografía una lámina.
\setbeamertemplate{frame footline}{%
  \begin{beamercolorbox}[wd=\textwidth, sep=3ex]{footline}%
    \usebeamerfont{page number in head/foot}%
    \usebeamercolor[fg]{page number in head/foot}%
    \nombreCursoCorto\hfill\usebeamertemplate*{frame numbering}%
  \end{beamercolorbox}%
}

%------------------- Notas del ponente (segunda pantalla) --------------
% Se activa desde el Makefile: `make notas` compila una versión con las
% notas al lado, para el portátil del ponente. La de proyectar no las lleva.
\ifdefined\connotas
  \setbeameroption{show notes on second screen=right}
\fi
